\clearpage\thispagestyle{empty}
\pdfbookmark[0]{Quellen zur Besprechung der Mailänder Ephemeriden}{effemeridi1813-sources}
\begin{center}\textbf{Quellen zur Besprechung der Mailänder Ephemeriden}\end{center}
\noindent KI-Übertragung des deutschen Originals von Carl Friedrich Gauß in modernes Deutsch nach den \href{https://gdz.sub.uni-goettingen.de/id/PPN236007467}{\emph{Werken}, Band VI}, Göttingen 1874, S.\ 556--558. Die Besprechung erschien zuerst unter Mailand in den \emph{Göttingischen gelehrten Anzeigen}, Stück 95, 14. Juni 1813, S.\ 945--952. In der Werkausgabe endet die Piazzi-Epoche mit einem Zeichen, das wie ein \emph{c} aussieht, und der Vergleichswert von 30 Sekunden gilt für heitere Sonnentage. Im Erstdruck stehen 1810 und heitere Sommertage. Beide Drucke geben für das Komplement im Abschnitt über Pallas 42 Stunden an.

\par\smallskip\noindent Francesco Carlini und Carlo Brioschi berechneten die Ephemeriden für \href{https://books.google.it/books?id=GDlZAAAAcAAJ}{1812} und \href{https://books.google.it/books?id=4HCOwKAJ28kC}{1813}. Die verlinkten Digitalisate stammen aus der Österreichischen Nationalbibliothek und der Biblioteca di Matematica G.~Castelnuovo der Sapienza Università di Roma. Der \href{https://www.beniculturali.inaf.it/opac/pst-iccd/circolo-moltiplicatore_13070}{Brera-Katalog zum Reichenbachschen Wiederholungskreis} verzeichnet dessen Tragsäule und verweist auf Barnaba Orianis Beobachtungen im Anhang für 1812.

\par\smallskip\noindent Erstdruck, digitalisiert von der Niedersächsischen Staats- und Universitätsbibliothek Göttingen, \href{https://images.p.sub.uni-goettingen.de/iiif/image/gdz:PPN319721507_1813:00000949/full/1650,/0/default.jpg}{S.\ 945}, \href{https://images.p.sub.uni-goettingen.de/iiif/image/gdz:PPN319721507_1813:00000950/full/1650,/0/default.jpg}{S.\ 946}, \href{https://images.p.sub.uni-goettingen.de/iiif/image/gdz:PPN319721507_1813:00000951/full/1650,/0/default.jpg}{S.\ 947}, \href{https://images.p.sub.uni-goettingen.de/iiif/image/gdz:PPN319721507_1813:00000952/full/1650,/0/default.jpg}{S.\ 948}, \href{https://images.p.sub.uni-goettingen.de/iiif/image/gdz:PPN319721507_1813:00000953/full/1650,/0/default.jpg}{S.\ 949}, \href{https://images.p.sub.uni-goettingen.de/iiif/image/gdz:PPN319721507_1813:00000954/full/1650,/0/default.jpg}{S.\ 950}, \href{https://images.p.sub.uni-goettingen.de/iiif/image/gdz:PPN319721507_1813:00000955/full/1650,/0/default.jpg}{S.\ 951} und \href{https://images.p.sub.uni-goettingen.de/iiif/image/gdz:PPN319721507_1813:00000956/full/1650,/0/default.jpg}{S.\ 952}.

\par\smallskip\noindent Übertragung und Satz wurden von KI erstellt.
\clearpage
